\documentclass[a4paper]{article}
\usepackage[affil-it]{authblk}
\usepackage{cite}
\usepackage{amsmath,amssymb,amsfonts,amsthm}
\usepackage{geometry}
\usepackage{booktabs}
\geometry{a4paper,
  left=1.5cm,
  right=1.5cm,
  top=1.5cm,  
  bottom=1.5cm,
  headheight=12pt,
  headsep=8pt}

\usepackage{enumitem}

\bibliographystyle{plain}

\begin{document}
% =================================================
\title{Numerical Analysis Homework 1}
\author{Zhuqingyu 3230105823}
\thanks{Electronic address: \texttt{3230105823@zju.edu.cn}}
\affil{(Mathematics and Applied Mathematics 2301), Zhejiang University}
\date{Due time: \today}
\maketitle
% ============================================
\section*{Theoretical Questions}

\subsection*{Question 1}
Consider the bisection method starting with the initial interval $[1.5,3.5]$. In the following questions ``the interval'' refers to the bisection interval whose width changes across different loops.
\begin{enumerate}[label=(\alph*)]
    \item What is the width of the interval at the $n$th step?
    \item What is the supremum of the distance between the root $r$ and the midpoint of the interval?
\end{enumerate}

\textbf{Solution:}
\begin{enumerate}[label=(\alph*)]
    \item Since $h_{n} = \frac{h_{n-1}}{2} = \frac{h_0}{2^{n}} = \frac{b_0 - a_0}{2^{n}} = \frac{3.5-1.5}{2^{n}} = \frac{1}{2^{n-1}}$.
    
    \item Since the root $\alpha \in [a_n, b_n]$ and $b_n - a_n = h_n = \frac{1}{2^{n-1}}$, hence $\lvert c_n - \alpha \rvert \leq \frac{1}{2}(b_n - a_n) = \frac{1}{2^{n}}$.
\end{enumerate}

\subsection*{Question 2}
In using the bisection algorithm with its initial interval as $[a_0,b_0]$ with $a_0>0$, we want to determine the root with its \textit{relative} error no greater than $\epsilon$. Prove that this goal of accuracy is guaranteed by the following choice of the number of steps,
\[
n \geq \frac{\log(b_0 - a_0) - \log\epsilon - \log a_0}{\log 2} - 1.
\]

\textbf{Solution:}
To make sure that $\frac{\lvert c_n -\alpha\rvert}{\alpha}<\epsilon$, since
$\frac{\lvert c_n-\alpha \rvert}{\lvert\alpha\rvert}\leq\frac{\lvert c_n-\alpha \rvert}{\lvert a_n\rvert}\leq\frac{\lvert c_n-\alpha \rvert}{\lvert a_0 \rvert}\leq\frac{b_0-a_0}{2^{n+1}a_0}<\epsilon$, which implies that 
\[
n \geq \frac{\log(b_0 - a_0) - \log \epsilon - \log a_0}{\log 2} - 1.
\]

\subsection*{Question 3}
Perform four iterations of Newton's method for the polynomial equation $p(x) = 4x^3 - 2x^2 + 3 = 0$ with the starting point $x_0 = -1$. Use a hand calculator and organize results of the iterations in a table.

\textbf{Solution:}
$p(x) = 4x^3 - 2x^2 + 3$ has the derivative $p'(x) = 12x^2 - 4x$.
use the initial guess $x_0 = -1$ and $x_{n+1}=x_n - \frac{p(x_n)}{p'(x_n)}$, we have

\begin{table}[ht]
\centering
\caption{Newton's Method Iterations}
\begin{tabular}{ccccc}
\toprule
Iteration & $x_n$ & $p(x_n)$ & $p'(x_n)$ & $x_{n+1}$ \\
\midrule
0 & -1.000000 & -3.000000 & 16.000000 & -0.812500 \\
1 & -0.812500 & -0.465820 & 11.171875 & -0.770804 \\
2 & -0.770804 & -0.020138 & 10.212886 & -0.768832 \\
3 & -0.768832 & -0.000043 & 10.168568 & -0.768828 \\
\bottomrule
\end{tabular}
\end{table}

\subsection*{Question 4}
Consider a variation of Newton's method in which only the derivative at $x_0$ is used,
\[
x_{n+1} = x_n - \frac{f(x_n)}{f'(x_0)}.
\]
Find $C$ and $s$ such that
\[
e_{n+1} = C e_n^s,
\]
where $e_n$ is the error of Newton's method at step $n$, $s$ is a constant, and $C$ may depend on $x_n$, the true solution $\alpha$, and the derivative of the function $f$.

\textbf{Solution:}
Let $g(x)=x-\frac{f(x)}{f'(x_0)}$, then we have $g'(x)=1-\frac{f'(x)}{f'(x_0)}$ and $g(\alpha)=\alpha$, so when $x_0$ is close to $\alpha$ (which is a root of $f$), i.e.
$\exists \delta>0, \forall x\in[\alpha-\delta,\alpha+\delta], f'(x)\ne0$, take $x_0 \in [\alpha-\delta,\alpha+\delta]$
\[
x_{n+1}-\alpha=g(x_n)-g(\alpha)=\left(1-\frac{f'(\eta_n)}{f'(x_0)}\right)(x_n-\alpha),
\]
where $\eta_n$ is between $x_n$ and $\alpha$. Thus $C = 1-\frac{f'(\eta_n)}{f'(x_0)}$ and $s = 1$.

\subsection*{Question 5}
Within $(-\frac{\pi}{2}, \frac{\pi}{2})$, will the iteration $x_{n+1} = \tan^{-1} x_n$ converge?

\textbf{Solution:}
The iteration $x_{n+1} = \arctan(x_n)$ converges:

Define $h(x)=\arctan x-x$, $h'(x)=\frac{1}{1+x^2}-1<0,\forall x \in (-\pi/2,\pi/2)$, so $h(x)$ is decreasing.
Since $h(0)=0$, hence when $x>0$, $h(x)<0$, when $x<0$, $h(x)>0$. If $x_0>0$, by induction we have $x_{n+1}-x_n=\arctan(x_n)-x_n<0$, so $0<x_{n+1}<x_n$. If $x_0<0$, by induction we have $x_{n+1}-x_n=\arctan(x_n)-x_n>0$, so $x_n<x_{n+1}<0$. Using monotonic sequence theorem, $x_n$ converges.

Assume that $x_n$ converges to $\beta$, since $\arctan(x)$ is continuous, we have $\beta = \arctan(\beta)$, so $\beta=0$, $x_n \to 0$ as $n \to \infty$.

\subsection*{Question 6}
Let $p > 1$. What is the value of the following continued fraction?
\[
x = \dfrac{1}{p + \dfrac{1}{p + \dfrac{1}{p + \cdots}}}
\]
Prove that the sequence of values converges. (Hint: this can be interpreted as $x = \lim_{n \to \infty} x_n$, where $x_1 = \frac{1}{p}$, $x_2 = \frac{1}{p + \frac{1}{p}}$, $x_3 = \frac{1}{p + \frac{1}{p + \frac{1}{p}}}$, \ldots, and so forth. Formulate $x$ as a fixed point of some function.)

\textbf{Solution:}
$x_{n+1}=\frac{1}{p+x_n}$ and $x_n \in (0,1/p)$. Define $g(x)=\frac{1}{p+x}$, $g'(x)=\frac{-1}{(p+x)^2}$, $\lvert g'(x)\rvert<1$ for all $x\in(0,1/p)$, so $g(x)$ is a contraction mapping.
Thus $x_n \to \alpha$ as $n \to \infty$, where $\alpha$ is a fixed point of $g$. Solving $\alpha = \frac{1}{p+\alpha}$ gives $\alpha = \frac{-p+\sqrt{p^2+4}}{2}$.

\subsection*{Question 7}
What happens in problem 2 if $a_0 < 0 < b_0$? Derive an inequality of the number of steps similar to that in problem 2. In this case, is the relative error still an appropriate measure?

\textbf{Solution:}
Consider example $y=\arctan x-\delta$, $\alpha_{\delta}=\tan(\delta)$, take small $\delta$ or let $\delta =0$, relative error $\frac{\lvert x_{n}-\alpha_{\delta}\rvert}{\alpha_{\delta}}$ can't measure the accuracy of the iteration, so we need to use absolute error $\lvert x_{n}-\alpha_{\delta}\rvert$.
$\lvert x_{n}-\alpha\rvert \leq \frac{1}{2^{n+1}}(b_0 - a_0) \leq \epsilon$ $\implies$
$n>\frac{log(b_0 - a_0) - \log \epsilon}{log 2}-1$
\subsection*{Question 8}
($*$) Consider solving $f(x) = 0$ ($f \in \mathcal{C}^{k+1}$) by Newton's method with the starting point $x_0$ close to a root of multiplicity $k$. Note that $\alpha$ is a zero of multiplicity $k$ of the function $f$ iff
\[
f^{(k)}(\alpha) \neq 0; \qquad \forall i < k, \quad f^{(i)}(\alpha) = 0.
\]
\begin{enumerate}[label=(\alph*)]
    \item How can a multiple zero be detected by examining the behavior of the points $(x_n, f(x_n))$?
    \item Prove that if $r$ is a zero of multiplicity $k$ of the function $f$, then quadratic convergence in Newton's iteration will be restored by making this modification:
    \[
    x_{n+1} = x_n - k \frac{f(x_n)}{f'(x_n)}.
    \]
\end{enumerate}

\textbf{Solution:}
\begin{enumerate}[label=(\alph*)]
    \item A multiple zero can be detected by observing that the convergence of Newton's method becomes linear rather than quadratic when approaching a multiple root.
    
    \item \begin{proof}
    Since $f \in C^{k+1}$ and $\alpha$ is a root of multiplicity $k$, we have $f' \in C^k$, $f'' \in C^{k-1}$. Then:
    
    \begin{align*}
    f(x) &= (x-\alpha)^k \cdot \frac{f^{(k)}(\alpha)}{k!} + O\left((x-\alpha)^{k+1}\right) \\
    f'(x) &= \frac{f^{(k)}(\alpha)}{(k-1)!} (x-\alpha)^{k-1} + O\left((x-\alpha)^k\right) \\
    f''(x) &= \frac{f^{(k)}(\alpha)}{(k-2)!} (x-\alpha)^{k-2} + O\left((x-\alpha)^{k-1}\right)
    \end{align*}
    
    We prove that Newton's method has linear convergence. Define:
    \[
    h(x) = x - \frac{f(x)}{f'(x)}
    \]
    Then:
    \[
    h'(x) = 1 - \frac{1}{(f'(x))^2} \cdot \left[(f'(x))^2 - f(x) \cdot f''(x)\right] = \frac{f(x) f''(x)}{(f'(x))^2}
    \]
    
    Note that $h'(x)$ is not defined at $x = \alpha$ since $f'(\alpha) = 0 = f(\alpha)$. However, we can compute the limit as $x \to \alpha$:
    
    \begin{align*}
    h'(x) &= \frac{\left((x-\alpha)^k \cdot \frac{f^{(k)}(\alpha)}{k!} + o((x-\alpha)^k)\right) \cdot \left(\frac{f^{(k)}(\alpha)}{(k-2)!} (x-\alpha)^{k-2} + o((x-\alpha)^{k-2})\right)}{\left(\frac{f^{(k)}(\alpha)}{(k-1)!} (x-\alpha)^{k-1} + o((x-\alpha)^{k-1})\right)^2} \\
    &= \frac{\frac{f^{(k)}(\alpha)^2}{k!(k-2)!} (x-\alpha)^{2k-2} + o((x-\alpha)^{2k-2})}{\frac{f^{(k)}(\alpha)^2}{((k-1)!)^2} (x-\alpha)^{2k-2} + o((x-\alpha)^{2k-2})} \\
    &\to \frac{(k-1)!^2}{k!(k-2)!} = \frac{((k-1)!)^2}{k!(k-2)!} = \frac{k-1}{k} = 1 - \frac{1}{k} \quad (k \geq 2)
    \end{align*}
    
    So we define $h'(\alpha) = 1 - \frac{1}{k}$ and $h \in C^1$ near $\alpha$, with $|h'(\alpha)| < 1$. Therefore, Newton's method has linear convergence, and the sequence $(x_n, f(x_n))$ will converge slower than expected.
    
    Define the modified iteration function:
    \[
    r(x) = x - k \cdot \frac{f(x)}{f'(x)} = k\left(x - \frac{f(x)}{f'(x)}\right) - (k-1)x = k h(x) - (k-1)x
    \]
    Then:
    \[
    r'(\alpha) = k \cdot \left(1 - \frac{1}{k}\right) - (k-1) = (k-1) - (k-1) = 0
    \]
    
    Using Corollary 1.42\cite{ZQH_NPDE2025}, we know that at least quadratic convergence will be restored.
    \end{proof}
\end{enumerate}

% ===============================================
\bibliography{Chapter1}
\end{document}